% =====================================================================
%  Preface / 前言
%  Gattringer & Lang, "Quantum Chromodynamics on the Lattice" (LNP 788)
%  Reconstructed from extract/preface.txt.
% =====================================================================

\chapter*{前言}
\addcontentsline{toc}{chapter}{前言}
\markboth{前言}{前言}

量子色动力学（QCD）是描述夸克与胶子的基本量子场论。为了在数学上严格地讨论它，该理论必须加以正则化。将时空替换为欧几里得格点已被证明是一种有效的方法，它既便于理论理解，又便于计算分析。格点 QCD 已经成为基本粒子物理中的一种标准工具。

正如书名所示：这是一本入门书！本书面向该领域的新手，作为他们学习的起点。我们只是想要一本能够放入高年级学生手中的书，供他们初读格点 QCD。这位想象中的学生需要具备高等量子力学、一定连续量子场论以及基本粒子物理唯象学基础知识作为前提。

鉴于当前研究中应用之广泛，这里所呈现的主题是有限的，我们不得不做出一些痛苦的选择。我们讨论 QCD，但省略了大多数其他的格点场论应用，例如标量理论、规范–Higgs 模型或电弱理论。尽管我们试图引导读者达到当今的理解水平，但我们不可能涵盖所有正在进行的活动，特别是涉及 QCD 在电弱理论中所起作用的部分。诸如胶球、拓扑激发等课题，以及手征微扰论等方法，仅被简要提及。这使我们能够相当详尽地覆盖其他主题，包括对关键方程的详细推导。该领域正在迅速发展。年度格点会议的会议文集提供了关于更新方向与最新结果的信息。

和往常一样，完成本书所花的时间比原计划要长，我们感谢编辑 Claus Ascheron 的耐心。我们非常感谢许多同事，他们愿意审读其中的这一篇或那一篇。我们特别要感谢 Vladimir Braun、Dirk Br\"ommel、Tommy Burch、Stefano Capitani、Tom DeGrand、Stephan D\"urr、Georg Engel、Christian Hagen、Leonid Glozman、Meinulf G\"ockeler、Peter Hasenfratz、Jochen Heitger、Verena Hermann、Edwin Laermann、Markus Limmer、Pushan Majumdar、Daniel Mohler、Wolfgang Ortner、Bernd-Jochen Schaefer、Stefan Schaefer、Andreas Sch\"afer、Erhard Seiler、Stefan Sint、Stefan Solbrig 和 Pierre van Baal。

如果本书中没有错误，那反倒令人惊讶了。因此我们为本书建立了一个网站配套页面：\url{http://physik.uni-graz.at/qcdlatt/}
在该页面上我们记录了勘误，并提供进一步的链接和信息。

\begin{flushright}
格拉茨，2009 年 3 月\\[2mm]
Christof Gattringer\\[1mm]
Christian B. Lang
\end{flushright}
